\chapter{DỰ ÁN CUỐI KHÓA: CHATBOT NHIỀU CUỘC HỘI THOẠI}
\label{chap:capstone}

\section{Mô tả đề tài và Mục tiêu học tập}

Dự án Capstone cuối khóa đóng vai trò tổng hợp toàn bộ các kỹ năng thiết kế hệ thống hội thoại có bộ nhớ. Học viên được yêu cầu nâng cấp toàn diện hệ thống cổng dịch vụ (LLM Gateway) viết bằng FastAPI kết hợp giao diện Streamlit nhằm đạt được các mục tiêu kỹ thuật thực tế: cho phép một người dùng tạo và quản lý song song nhiều phiên hội thoại độc lập, đảm bảo lịch sử trò chuyện được lưu trữ bền vững qua các lần khởi động lại máy chủ, kiểm soát chặt chẽ kích thước ngữ cảnh gửi đi, và chỉ lưu thông tin bộ nhớ dài hạn khi người dùng cho phép rõ ràng.

\section{Yêu cầu chức năng chi tiết}

Học viên phải hoàn thành tối thiểu các tính năng chức năng sau đây:

\begin{enumerate}
  \item \textbf{Quản lý phiên hội thoại (Thread Management):} Giao diện người dùng (Streamlit) hỗ trợ các tính năng tạo luồng chat mới, hiển thị danh sách các luồng chat cũ, cho phép đổi tên tiêu đề luồng chat và xóa vĩnh viễn luồng chat.
  \item \textbf{Tương tác không trạng thái phía Client:} Client gửi tin nhắn mới lên server chỉ truyền kèm \texttt{thread\_id} và \texttt{user\_id}; tuyệt đối không gửi kèm các tin nhắn cũ trong payload.
  \item \textbf{Khôi phục lịch sử sau khi làm mới giao diện:} Khi người dùng nhấn F5 làm mới trình duyệt, frontend tự động truy vấn endpoint lấy lịch sử tin nhắn để vẽ lại toàn bộ các bong bóng chat cũ một cách chính xác.
  \item \textbf{Không trộn lẫn trạng thái giữa các Nhà cung cấp (Providers):} Cho phép người dùng chuyển đổi qua lại giữa DeepSeek và Gemini tại giao diện cấu hình runtime mà không làm mất hoặc trộn lẫn lịch sử tin nhắn của luồng hiện tại.
  \item \textbf{Cắt bớt tin nhắn tự động (Auto-trimming):} Tích hợp cơ chế tự động tính toán token và cắt bớt tin nhắn đầu vào theo token budget trước khi gọi LLM API.
  \item \textbf{Tóm tắt lũy tiến (Running Summary):} Tự động kích hoạt node tóm tắt lịch sử cũ khi tổng kích thước tin nhắn vượt quá ngưỡng token cấu hình.
  \item \textbf{Ghi nhớ chủ động (Opt-in Long-term Memory):} Trợ lý chỉ được phép ghi nhận các tùy chọn cá nhân hoặc dữ kiện dài hạn vào Store khi nhận được câu lệnh bắt đầu bằng từ khóa ``Hãy nhớ \dots''. Trợ lý phải hỗ trợ câu lệnh ``Hãy quên \dots'' để thực hiện xóa bỏ thông tin tương ứng.
  \item \textbf{Bảo mật cách ly người dùng tuyệt đối:} Ngăn chặn triệt để hành vi rò rỉ dữ liệu chéo (cross-user data leak). Người dùng A tuyệt đối không được phép đọc, ghi hoặc xóa luồng chat hay bộ nhớ Store thuộc sở hữu của người dùng B.
\end{enumerate}

\section{Yêu cầu phi chức năng và Vận hành}

\begin{itemize}
  \item \textbf{An toàn thông tin (Security):} Tuyệt đối không lưu các khóa bí mật (API Keys), mật khẩu CSDL hoặc tokens xác thực dưới dạng văn bản thô trong mã nguồn hoặc log. Cấu hình tất cả thông qua biến môi trường.
  \item \textbf{Xử lý ngoại lệ (Exception Handling):} Thiết lập cơ chế timeout (ví dụ: 15 giây) và tự động thử lại (retry) khi gọi LLM API. Trả về thông báo lỗi có nghĩa ngoài giao diện thay vì làm crash tiến trình backend.
  \item \textbf{Độ bền dữ liệu (Durability):} Tích hợp SQLite hoặc PostgreSQL làm checkpointer vật lý lưu trữ checkpoint, chứng minh dữ liệu khôi phục hoàn chỉnh sau khi khởi động lại backend.
  \item \textbf{Giám sát hiệu năng (Observability):} In ra màn hình log hoặc lưu trữ metric về số lượng token tiêu thụ cho mỗi request và thời gian phản hồi (latency) của từng node trong đồ thị.
  \item \textbf{Tài liệu kỹ thuật (Documentation):} Viết tệp README mô tả chi tiết luồng di chuyển của dữ liệu (data flow) từ Streamlit qua FastAPI đến LangGraph, sơ đồ database, và chính sách thời hạn lưu giữ dữ liệu (retention policy).
\end{itemize}

\section{Kiến trúc tham chiếu hệ thống Capstone}

Sơ đồ dưới đây mô tả sự tương tác giữa các thành phần trong kiến trúc tham chiếu mà học viên cần triển khai:

\begin{figure}[H]
\centering
\resizebox{0.96\textwidth}{!}{%
\begin{tikzpicture}[
 component/.style={draw=aiBlue,rounded corners,fill=aiSoftBlue,minimum width=3.2cm,minimum height=1cm,align=center},
 data/.style={draw=aiGold,rounded corners,fill=aiSoftGold,minimum width=3.2cm,minimum height=1cm,align=center},
 arr/.style={-{Stealth[length=2mm]},thick,aiNavy}]
 \node[component] (ui) {Streamlit\\Giao diện người dùng};
 \node[component,right=1.0cm of ui] (api) {FastAPI routes\\Phân quyền \& Xác thực};
 \node[component,right=1.0cm of api] (service) {ChatService\\Điều phối \& Chính sách};
 \node[component,below=1.1cm of service] (graph) {LangGraph\\Đồ thị điều khiển};
 \node[data,left=1.0cm of graph] (cp) {Database\\Checkpoint Saver};
 \node[data,right=1.0cm of graph] (store) {Database Store\\User Memory};
 \draw[arr] (ui)--(api); \draw[arr] (api)--(service); \draw[arr] (service)--(graph);
 \draw[arr] (graph)--(cp); \draw[arr] (graph)--(store);
\end{tikzpicture}
}
\caption{Kiến trúc tham chiếu cho dự án cuối khóa.}
\label{fig:capstone-architecture}
\figsource{Tác giả tự dựng.}
\end{figure}

\section{Kế hoạch thực hiện theo mốc kiểm chứng}

\begin{longtable}{p{1.2cm}p{5.2cm}p{8.2cm}}
\caption{Các mốc dự án có thể kiểm chứng}\label{tab:capstone-slices}\\
\toprule
\textbf{Mốc} & \textbf{Phạm vi kiểm thử} & \textbf{Điều kiện nghiệm thu (Definition of Ready)} \\
\midrule
\endfirsthead
\toprule
\textbf{Mốc} & \textbf{Phạm vi kiểm thử} & \textbf{Điều kiện nghiệm thu (Definition of Ready)} \\
\midrule
\endhead
1 & API Conversation & Tạo thành công hai luồng chat bằng mô hình giả lập thông qua Swagger; kiểm thử phân quyền 403 hoạt động chính xác. \\
2 & Adapter gọi mô hình thực & Kết nối thành công mô hình DeepSeek/Gemini; xử lý tốt lỗi mất kết nối mạng hoặc lỗi định dạng tin nhắn. \\
3 & Lưu trữ vật lý bền vững & Thay thế saver bằng SQLite/PostgreSQL; restart server và kiểm chứng lịch sử không bị mất. \\
4 & Tối ưu hóa Context & Cấu hình trim tin nhắn theo token budget; kiểm thử thành công kịch bản chat siêu dài (trên 50 tin nhắn). \\
5 & Bộ nhớ dài hạn Store & Người dùng nói ``Hãy nhớ tôi thích Python''; kiểm tra tài khoản đó được nhớ ở luồng chat mới; kiểm tra tài khoản khác không thấy thông tin này. \\
6 & Hoàn thiện UI và Release & Giao diện Streamlit vẽ mượt mà luồng trò chuyện, hiển thị thông tin đo đạc token và thời gian phản hồi ở mỗi câu trả lời. \\
\bottomrule
\end{longtable}

\section{Kịch bản Demo 5 phút bắt buộc}

Khi nghiệm thu dự án, học viên phải thực hiện trình diễn trực tiếp theo đúng 6 bước kịch bản sau:
\begin{enumerate}
  \item \textbf{Bước 1 (Khai báo thông tin):} Tại luồng hội thoại A, nói với trợ lý ``Tôi tên Nhật''.
  \item \textbf{Bước 2 (Kiểm tra cách ly luồng ngắn hạn):} Khởi tạo luồng hội thoại B và hỏi ``Tôi tên gì?''. Trợ lý phải trả lời chưa biết tên (chứng minh short-term memory cách ly).
  \item \textbf{Bước 3 (Ghi nhớ dài hạn chủ động):} Quay lại luồng hội thoại A, nói ``Hãy nhớ tôi tên Nhật''. Trợ lý phản hồi xác nhận sẽ ghi nhớ.
  \item \textbf{Bước 4 (Kiểm tra chia sẻ bộ nhớ dài hạn):} Khởi tạo luồng hội thoại C mới hoàn toàn của cùng người dùng đó, thực hiện câu hỏi ``Tôi tên gì?''. Trợ lý phải trả lời tên Nhật (chứng minh bộ nhớ dài hạn Store hoạt động xuyên luồng).
  \item \textbf{Bước 5 (Kiểm tra cách ly tài khoản người dùng):} Đăng nhập bằng tài khoản người dùng B (hoặc gửi request bằng API Client giả lập người dùng B) và thực hiện câu hỏi hỏi tên. Trợ lý tuyệt đối không được phép trả lời tên Nhật.
  \item \textbf{Bước 6 (Kiểm tra độ bền lưu trữ):} Thực hiện lệnh restart dịch vụ backend FastAPI. F5 lại giao diện người dùng Streamlit và chứng minh lịch sử trò chuyện vẫn được nạp và hiển thị đầy đủ.
\end{enumerate}

\section{Định nghĩa hoàn thành (Definition of Done)}

Một dự án Capstone được tính là hoàn thành xuất sắc khi và chỉ khi:
\begin{itemize}
  \item Tất cả các mốc kiểm chứng từ 1 đến 6 hoạt động ổn định và vượt qua buổi trình diễn demo 5 phút trực tiếp.
  \item Bộ mã nguồn có độ bao phủ kiểm thử tự động (test coverage) đạt trên 80\%, vượt qua toàn bộ các bài test của pytest.
  \item Hệ thống không chứa bất kỳ mã khóa API hay mật khẩu thô nào trong mã nguồn.
  \item Có tệp tài liệu README chi tiết và chính sách quản lý thời gian lưu trữ dữ liệu rõ ràng.
\end{itemize}

\begin{aicheckpoint}[Tự phản biện thiết kế hệ thống]
Học viên cần tự viết một đoạn phản biện trả lời câu hỏi: Tại sao chúng ta nên sử dụng mô hình quản lý bộ nhớ của LangGraph thay vì tự viết các câu lệnh truy vấn SQL thủ công để lưu lịch sử vào database? Những phần việc nào của bộ nhớ vẫn thuộc về trách nhiệm xử lý của tầng ứng dụng? Nếu bỏ LangGraph, bạn sẽ gặp khó khăn gì khi triển khai các quy trình tác tử phức tạp hơn (ví dụ như Human-in-the-loop)?
\end{aicheckpoint}
